
\documentclass{article}
\usepackage{colortbl}
\usepackage{makecell}
\usepackage{multirow}
\usepackage{supertabular}

\begin{document}

\newcounter{utterance}

\centering \large Interaction Transcript for game `codenames', experiment `opponent\_difficult', episode 9 with reflexion{-}Qwen3.5{-}27B.
\vspace{24pt}

{ \footnotesize  \setcounter{utterance}{1}
\setlength{\tabcolsep}{0pt}
\begin{supertabular}{c@{$\;$}|p{.15\linewidth}@{}p{.15\linewidth}p{.15\linewidth}p{.15\linewidth}p{.15\linewidth}p{.15\linewidth}}
    \# & $\;$A & \multicolumn{4}{c}{Game Master} & $\;\:$B\\
    \hline

    \theutterance \stepcounter{utterance}  
    & & & \multicolumn{2}{p{0.3\linewidth}}{
        \cellcolor[rgb]{0.9,0.9,0.9}{
            \makecell[{{p{\linewidth}}}]{
                \texttt{\tiny{[GM$|$GM]}}
                {'hidden': {'team': ['plot'], 'innocent': ['sink', 'millionaire', 'tail', 'nurse', 'figure'], 'opponent': [], 'assassin': []}, 'revealed': {'team': {'team': [], 'innocent': [], 'opponent': [], 'assassin': []}, 'opponent': {'team': [], 'innocent': [], 'opponent': [], 'assassin': []}}}
            }
        }
    }
    & & \\ \\

    \theutterance \stepcounter{utterance}  
    & & \multicolumn{4}{p{0.6\linewidth}}{
        \cellcolor[rgb]{0.9,0.9,0.9}{
            \makecell[{{p{\linewidth}}}]{
                \texttt{\tiny{[P1$\langle$GM]}}
                \texttt{You have attempted this same task before and did not succeed. These are the plans you wrote after those attempts.} \\
\texttt{Trial \#1: New plan: My previous strategy failed because I prioritized grouping multiple words under a single clue (e.g., "ring" for Saturn, bar, and plot) without sufficiently verifying that the association was strong enough to avoid triggering distractors or ending the turn prematurely. Specifically, "ring" likely led to the distractor "figure" or "queen" (associated with rings) or simply failed to connect clearly to "plot" and "bar," causing the turn to end before all words were guessed. In the future, I will prioritize clarity over quantity: I will select clues that have a very direct, unambiguous link to 1 or 2 specific remaining words rather than trying to force a connection to 3 or more. If a word like "plot" or "bar" is difficult to group, I will target it individually with a precise clue (e.g., "story" for plot, "pub" for bar) to ensure the partner guesses it immediately, even if it means taking more turns, rather than risking a turn{-}ending mistake with a vague multi{-}word clue.} \\
\\ 
\texttt{Trial \#2: New plan: My previous strategy failed because I attempted to group all three words ("saturn", "bar", "plot") under a single clue ("RING") that was too ambiguous, causing my partner to guess a distractor ("queen") and end the turn prematurely. While my subsequent individual clues were correct, the initial loss of the turn prevented a full sweep. In the future, I will prioritize safety and turn preservation over maximizing words per turn. I will analyze the semantic distance between my target words and the opponent/distractor lists before grouping. If a single clue risks triggering a distractor (as "ring" did with "queen"), I will immediately split the targets into smaller, safer groups (e.g., 1 or 2 words) or target them individually to ensure the partner guesses correctly without interruption. I will only attempt multi{-}word clues if the association is exclusive to my team's words and clearly distinct from all distractors.} \\
\\ 
\texttt{Trial \#3: New plan: My previous attempt failed because I provided a clue ("STORY") that, while semantically correct for "plot," was too broad and likely triggered a distractor word (such as "figure" as in a story figure, or "nurse" in a medical story context) or simply failed to distinguish the target from the distractors clearly enough to prevent a wrong guess. Although I only had one word to guess, the lack of specificity in the clue relative to the distractor list caused the turn to end prematurely. In the future, even with a single target word, I will select a clue that is highly specific to the target's unique context and explicitly distinct from all distractors. For "plot," instead of the generic "STORY," I would choose a more precise term like "LAND" (referring to a plot of land) or "SCHEME" (referring to a plot to do something), ensuring the association is strong for the target but weak or non{-}existent for "sink," "millionaire," "tail," "nurse," and "figure." I will always cross{-}reference my chosen clue against the distractor list to ensure no semantic overlap exists before submitting.} \\
\\ 
\\ 
\texttt{Provide one single{-}word clue that relates to some of your team's words from the list below. You can choose to how many and to which words your clue relates to.} \\
\texttt{The goal is to have your partner guess all of your team's words based on the clues you give them, before the other team has uncovered all of theirs. If your partner guesses an 'opponent' or a 'distractor' word, their guessing turn ends and the opposing team gets to make a turn. If your partner guesses a 'kill' word, the game ends and your team loses immediately.} \\
\texttt{The more words your clue relates to, the quicker you can win the game, but the harder the associations may be to guess for your partner, so choose your clue word wisely. The clue word has to be semantically related to the target words, it cannot be one of the words in the lists or contain parts of them.} \\
\\ 
\texttt{Always give your single{-}word clue and your comma{-}separated list of related target words in the following format and make your answers as short as possible, never include any other text than is required in this form:} \\
\\ 
\texttt{CLUE: <WORD>} \\
\texttt{TARGETS: <TARGETS>} \\
\\ 
\texttt{Your team words are: plot.} \\
\texttt{Your opponent's team words are: .} \\
\texttt{Distractor words are: sink, millionaire, tail, nurse, figure.} \\
\texttt{The kill words are: .} \\
            }
        }
    }
    & \\ \\

    \theutterance \stepcounter{utterance}  
    & \multicolumn{4}{p{0.6\linewidth}}{
        \cellcolor[rgb]{0.9,0.9,0.9}{
            \makecell[{{p{\linewidth}}}]{
                \texttt{\tiny{[P1$\rangle$GM]}}
                \texttt{CLUE: LAND} \\
\texttt{TARGETS: plot} \\
            }
        }
    }
    & & \\ \\

    \theutterance \stepcounter{utterance}  
    & & & \multicolumn{2}{p{0.3\linewidth}}{
        \cellcolor[rgb]{0.9,0.9,0.9}{
            \makecell[{{p{\linewidth}}}]{
                \texttt{\tiny{[GM$|$GM]}}
                \texttt{land} \\
            }
        }
    }
    & & \\ \\

    \theutterance \stepcounter{utterance}  
    & & & \multicolumn{2}{p{0.3\linewidth}}{
        \cellcolor[rgb]{0.9,0.9,0.9}{
            \makecell[{{p{\linewidth}}}]{
                \texttt{\tiny{[GM$|$GM]}}
                ['plot']
            }
        }
    }
    & & \\ \\

    \theutterance \stepcounter{utterance}  
    & & & \multicolumn{2}{p{0.3\linewidth}}{
        \cellcolor[rgb]{0.9,0.9,0.9}{
            \makecell[{{p{\linewidth}}}]{
                \texttt{\tiny{[GM$|$GM]}}
                {'word': 'plot', 'assignment': 'team'}
            }
        }
    }
    & & \\ \\

    \theutterance \stepcounter{utterance}  
    & & & \multicolumn{2}{p{0.3\linewidth}}{
        \cellcolor[rgb]{0.9,0.9,0.9}{
            \makecell[{{p{\linewidth}}}]{
                \texttt{\tiny{[GM$|$GM]}}
                \texttt{opponent has won} \\
            }
        }
    }
    & & \\ \\

    \theutterance \stepcounter{utterance}  
    & \multicolumn{4}{p{0.6\linewidth}}{
        \cellcolor[rgb]{0.95,0.95,0.95}{
            \makecell[{{p{\linewidth}}}]{
                \texttt{\tiny{[P2$\rightarrow$P1]}}
                \texttt{Post{-}episode Reflexion (not sent to GM):} \\
\texttt{New plan: My previous attempt failed because the clue "LAND" was too ambiguous; while it relates to "plot" (as in a plot of land), it likely triggered the distractor "tail" (as in the tail of a landmass or a specific idiom) or simply failed to provide a strong enough unique association for the partner to guess immediately, causing the turn to end. In the future, I will select a clue that targets the most distinct and common definition of the word while explicitly avoiding any semantic overlap with the distractors ("sink", "millionaire", "tail", "nurse", "figure"). For "plot", instead of "LAND", I will use "SCHEME" or "STORY", which strongly evoke the narrative or conspiratorial meaning of "plot" and have no logical connection to the specific distractors provided. I will always verify that the chosen clue does not share a common phrase or association with any distractor before submitting.} \\
            }
        }
    }
    & & \\ \\

\end{supertabular}
}

\end{document}
